\begin{trapbox}
	\begin{description}[style=nextline, leftmargin=2.8em, labelsep=0.5em, itemsep=0.35em]
		\item[\textbf{\textcolor{CTrap}{易错点一（符号正负）：}}] 当$B$或$C$为钝角时，$\cos B$或$\cos C$为负，若把投影误当成普通线段长度并取绝对值，会导致$a$错误。
		\item[\textbf{\textcolor{CTrap}{正确理解（有向投影）：}}] 射影定理中的余弦直接参与运算，$\cos B$和$\cos C$可以是负值，反映的是沿目标边方向的有向投影，直接代入即可。
		\item[\textbf{\textcolor{CTrap}{错因分析（锐角惯性）：}}] 只记住了锐角三角形作高线图形，忽略了钝角三角形中投影可能反向的情况。
	\end{description}

	\medskip
	\begin{description}[style=nextline, leftmargin=2.8em, labelsep=0.5em, itemsep=0.35em]
		\item[\textbf{\textcolor{CTrap}{易错点二（字母混淆）：}}] 将$a = b\cos C + c\cos B$误写为$a = b\cos B + c\cos C$，导致计算错误。
		\item[\textbf{\textcolor{CTrap}{正确理解（角边对应）：}}] 注意公式右端出现的是与$a$相邻两角$B,C$的余弦，而不是$a$的对角$A$的余弦；求$a$时需要$b,c,\cos B,\cos C$，不是只知道夹角$A$的余弦。
		\item[\textbf{\textcolor{CTrap}{错因分析（死记硬背）：}}] 没有理解投影的几何对称性，机械记忆导致对应关系错乱。
	\end{description}

	\medskip
	\begin{description}[style=nextline, leftmargin=2.8em, labelsep=0.5em, itemsep=0.35em]
		\item[\textbf{\textcolor{CTrap}{易错点三（记忆不全）：}}] 只记住一个等式，不会轮换写出另外两个，解题时缺乏工具。
		\item[\textbf{\textcolor{CTrap}{正确理解（循环对称）：}}] 三个等式可通过轮换$A,B,C$与$a,b,c$整体记忆，类似正弦定理的对称性。
		\item[\textbf{\textcolor{CTrap}{错因分析（孤立记忆）：}}] 只背了第一个，没有利用轮换性主动推导其余。
	\end{description}
\end{trapbox}